\par\Needspace{4\baselineskip}
\begin{center}
\SourceLeaf{023b}\pdfbookmark[2]{5.}{article-005}5.
\end{center}
The angle $v$ is confined between $-180^0$ and $+180^0$ in the parabola and between $-(180^0-\psi)$ and $+(180^0-\psi)$ in the hyperbola. In the ellipse, however, it traverses the full circle in cycles that repeat indefinitely. We call it the \emph{true anomaly} of the moving body. Until now, nearly all astronomers have been accustomed to measure the true anomaly in the ellipse from aphelion instead of perihelion, contrary to the analogous practice for the parabola and hyperbola, which have no aphelion and where the reckoning therefore had to start from perihelion. In restoring the analogy among all kinds of conic sections,
\SourceLeaf{024}we hesitated all the less because the most recent French astronomers had already set us an example.

It is often convenient to alter the form of the expression $r = \frac{p}{1+e\cos v}$ slightly. The following forms in particular are worth noting.
\[
r = \frac{p}{1+e-2e\sin\frac{1}{2}v^2} = \frac{p}{1-e+2e\cos\frac{1}{2}v^2}
\]
\[
r = \frac{p}{(1+e)\cos\frac{1}{2}v^2+(1-e)\sin\frac{1}{2}v^2}\cdot
\]

In the parabola we therefore have $r = \frac{p}{2\cos\frac{1}{2}v^2}$. For the hyperbola it is particularly convenient to use $r = \frac{p\cos\psi}{2\cos\frac{1}{2}(v+\psi)\cos\frac{1}{2}(v-\psi)}\cdot$
